\chapter{Flattening Stratification of a Coherent Sheaf}
\label{chap:Picard_FlatteningStratification}
% archon:covers AlgebraicJacobian/Picard/FlatteningStratification.lean

% STRATEGY NOTE (iter-174). This chapter is the A.2.a sub-row of Route~A
% (positive-genus arm of nonempty_jacobianWitness). It packages the
% Nitsure~\S4 flattening-stratification theorem, the technical heart that
% turns the absolute Quot-functor representability problem on a projective
% morphism \(X \to S\) with coherent input \(\F\) into a finite collection of
% locally-closed strata over which \(\F\) becomes flat. Mathlib (master
% \texttt{b80f227}) has neither a coherent-sheaf flatness predicate over a
% scheme base nor the flattening-stratification construction; this chapter
% is the prover-ready specification for the new file
% \texttt{AlgebraicJacobian/Picard/FlatteningStratification.lean} which a
% subsequent iter (file-skeleton lane) will scaffold.
%
% This chapter is independently parallel-startable: it does not gate on the
% Route~A.1.a relative-Spec chapter \cref{chap:Picard_RelativeSpec}.

\section{Setup and motivation}
\label{sec:flatstrat_setup}

The construction of the Quot scheme \(\Quot_{\F/X/S}^{\Phi,L}\) in Nitsure's exposition
of Grothendieck's existence theorem (\cite{nitsure-hilbert-quot}, \S5) proceeds by
embedding the functor of points of \(\Quot\) into a Grassmannian via Castelnuovo--Mumford
regularity, then cutting out the locus of \emph{quotients with flat fibres of the
prescribed Hilbert polynomial} by a closed-subscheme condition.  Cutting out this
locus uses, as a primary technical input, the existence of the
\emph{flattening stratification} of a coherent sheaf \(\F\) on a projective
\(S\)-scheme: a finite collection of locally-closed subschemes of \(S\) on each of which
the restriction of \(\F\) becomes flat.

This chapter packages the Nitsure~\S4 flattening-stratification theorem as the
foundational A.2.a building block. The Quot-scheme construction itself (A.2.b),
the relative Picard functor (A.1.c), and the assembly of the Picard scheme
representability theorem on a smooth proper curve (A.2.c) live in separate sibling
chapters.

We work throughout with a noetherian base scheme \(S\), a projective morphism
\(\pi : X \to S\), and a coherent \(\OO_X\)-module \(\F\). The conclusion of Nitsure's
theorem covers an arbitrary noetherian base with a projective \(X \to S\), which
specialises directly to the Route~A consumer setting where \(S = T\) is an arbitrary
\(k\)-scheme, \(X = C \times_k T\) for \(C\) a smooth proper curve over the base field
\(k\), and \(\F\) is a coherent quotient sheaf produced by the Grassmannian embedding
of the Quot functor (cf.\ \ below).

\section{Coherent flatness over a base}
\label{sec:coh_flat}

\begin{definition}

\leanok
[Coherent sheaf flat over the base]
  \label{def:coherent_sheaf_flat}
  \lean{AlgebraicGeometry.Scheme.CoherentSheafFlat}
  Let \(S\) be a noetherian scheme, \(f : X \to S\) a morphism of finite type, and
  \(\F\) a coherent \(\OO_X\)-module.  We say \(\F\) is \emph{flat over \(S\)}
  (equivalently, \(\F\) is \emph{\(f\)-flat}, or \emph{\(S\)-flat}) if for every
  point \(x \in X\) the stalk \(\F_x\) is a flat \(\OO_{S, f(x)}\)-module, where the
  \(\OO_{S, f(x)}\)-module structure is given by the composite
  \(\OO_{S, f(x)} \to \OO_{X, x} \to \End(\F_x)\).
  \textit{Source: cf.\ [Stacks Project], tag 00HB (flat modules); [Nitsure], \S4
  passim.}
  Equivalently (after pulling back to an affine open \(U = \Spec A \subseteq S\) and
  an affine open \(V = \Spec B \subseteq f^{-1}(U)\) with \(\F|_V = \widetilde{M}\)
  for \(M\) a finite \(B\)-module), \(\F|_V\) is \(\OO_U\)-flat in the sense that \(M\) is a
  flat \(A\)-module (with the \(A\)-module structure coming from the ring map
  \(A \to B\)).
\end{definition}

\begin{remark}
\label{rem:coh_flat_pullback}
  The condition is stable under base change: for any morphism \(g : S' \to S\),
  pulling back along the cartesian square
  \[
    \begin{tikzcd}
      X' \ar[r, "g'"] \ar[d, "f'"'] & X \ar[d, "f"] \\
      S' \ar[r, "g"]                & S
    \end{tikzcd}
  \]
  yields a coherent \(\OO_{X'}\)-module \(\F' := (g')^* \F\) which is \(S'\)-flat
  whenever \(\F\) is \(S\)-flat.  This is the routine cancellation of tensor-product
  associativity at the level of stalks. The converse fails in general.
\end{remark}

\section{Generic flatness}
\label{sec:generic_flatness}

The construction of the flattening stratification proceeds by Noetherian
induction on \(S\). The induction step relies on the following classical
\emph{generic flatness} theorem of Grothendieck: over a noetherian integral
base, a coherent sheaf is flat over a non-empty open subscheme. Algebraically,
this is the assertion that a finitely generated module over a finite-type
algebra over a noetherian domain is generically free after inverting a single
non-zero element of the base ring.

\begin{theorem}
[Generic flatness, algebraic form]
  \label{thm:generic_flatness_algebraic}
  \lean{AlgebraicGeometry.TODO.genericFlatnessAlgebraic}
  % SOURCE: [Nitsure], \S4 ``Lemma on Generic Flatness''
  % (read from references/nitsure-hilbert-quot-src/nitsure-hilbert-quot.tex, L1711-L1716).
  % SOURCE QUOTE:
  % "Let \(A\) be a noetherian domain,
  %  and \(B\) a finite type \(A\) algebra. Let \(M\) be a finite \(B\)-module.
  %  Then there exists an \(f\in A\), \(f\ne 0\), such that
  %  the localisation \(M_f\) is a free module over \(A_f\)."
  \textit{Source: [Nitsure], \S4, ``Lemma on Generic Flatness''.}
  Let \(A\) be a noetherian domain, \(B\) a finite-type \(A\)-algebra, and \(M\) a finite
  \(B\)-module. Then there exists \(f \in A\), \(f \ne 0\), such that the localisation
  \(M_f\) is a free \(A_f\)-module.
\end{theorem}

% SOURCE QUOTE PROOF: "Over the quotient field \(K\) of \(A\), the \(K\)-algebra
% \(B_K = K\otimes _AB\) is of finite type, and \(M_K = K\otimes_AM\) is a
% finite module over \(B_K\). Let \(n\) be the dimension of the support
% of \(M_K\) over \(\spec B_K\). We argue by induction on \(n\), starting
% with \(n=-1\) which is the case when \(M_K =0\). [...]
% Now let \(n\ge 0\), and let the lemma be proved for smaller values.
% As \(B\) is noetherian and \(M\) is assumed to be a finite \(B\)-module,
% there exists a finite filtration
% \(0=M_0 \subset \ldots  \subset M_r =M\)
% where each \(M_i\) is a \(B\)-submodule of \(M\) such that for each \(i\ge 1\)
% the quotient module \(M_i/M_{i-1}\) is isomorphic to \(B/\pp_i\) for
% some prime ideal \(\pp_i\) in \(B\). [...] By Noether normalisation lemma,
% there exist elements \(b_1,\ldots,b_n \in B\), such that \(K\otimes_AB\) is
% finite over its subalgebra \(K[b_1,\ldots, b_n]\) [...]"
% (Source: nitsure-hilbert-quot-src/nitsure-hilbert-quot.tex, L1719-L1772; full
% proof omitted for length, structural steps restated below in project notation.)
\begin{proof}
  
  Following [Nitsure], \S4: pass to the fraction field \(K = \mathrm{Frac}(A)\) and
  set \(B_K = K \otimes_A B\), \(M_K = K \otimes_A M\), \(n := \dim \mathrm{supp}_{\Spec B_K} M_K\).
  Argue by induction on \(n\), with base case \(n = -1\) (i.e.\ \(M_K = 0\)): each
  generator of \(M\) is annihilated by some non-zero element of \(A\); taking the
  product of finitely many annihilators yields \(f \in A\) with \(f M = 0\), hence
  \(M_f = 0\) is trivially free.  For \(n \ge 0\), a prime filtration of \(M\) as a
  finite \(B\)-module
  \[
    0 = M_0 \subset M_1 \subset \cdots \subset M_r = M
  \]
  with successive quotients \(M_i / M_{i-1} \cong B / \mathfrak{p}_i\) reduces the
  problem (via the splicing fact: \(M_{f'f''}\) is \(A_{f'f''}\)-free when \(M'_{f'}\)
  and \(M''_{f''}\) are free) to the case \(M = B/\mathfrak{p}\) for a prime
  \(\mathfrak{p} \subset B\), then to the case \(B\) a domain and \(M = B\).
  Noether normalisation applied to \(B_K / K\) yields algebraically independent
  \(b_1, \dots, b_n \in B\) with \(B_K\) finite over \(K[b_1, \dots, b_n]\); choosing
  \(g \ne 0\) to clear denominators of integrality equations, \(B_g\) is finite over
  \(A_g[b_1, \dots, b_n]\). Let \(m\) be the generic rank: the short exact sequence
  \[
    0 \to A_g[b_1, \dots, b_n]^{\oplus m} \to B_g \to T \to 0
  \]
  has \(T\) a finite torsion module over \(A_g[b_1, \dots, b_n]\), so $\dim
  \mathrm{supp}_{B_K} (K \otimes_{A_g} T) < n$.  The inductive hypothesis
  yields \(h \ne 0\) in \(A\) with \(T_h\) free over \(A_{gh}\); set \(f := gh\).
\end{proof}

\begin{theorem}

\leanok
[Generic flatness, geometric form]
  \label{thm:generic_flatness}
  \lean{AlgebraicGeometry.genericFlatness}
  \uses{thm:generic_flatness_algebraic}
  % NOTE (iter-303, review): the Lean signature `AlgebraicGeometry.genericFlatness`
  % takes `F : X.Modules` with NO coherence / finite-type hypothesis, so it is
  % FALSE as stated (the blueprint correctly requires `\F` coherent and `p`
  % finite-type). Its `sorry` is therefore unprovable. The decl is NOT in
  % archon-protected.yaml, so the planner must re-sign it to add the coherent /
  % LocallyOfFiniteType hypotheses this statement requires before any prover is
  % sent at it. See .archon/task_results/Picard_FlatteningStratification.md.

  % SOURCE: [Nitsure], \S4, Theorem ``generic flatness''
  % (read from references/nitsure-hilbert-quot-src/nitsure-hilbert-quot.tex, L1781-L1787).
  % SOURCE QUOTE:
  % "Let \(S\) be a noetherian and integral scheme.
  %  Let \(p: X\to S\) be a finite type morphism, and let \(\F\) be a
  %  coherent sheaf of \({\mathcal O}_X\)-modules. Then there exists a non-empty
  %  open subscheme \(U\subset S\) such that the restriction of \(\F\)
  %  to \(X_U = p^{-1}(U)\) is flat over \({\mathcal O}_U\)."
  \textit{Source: [Nitsure], \S4, Theorem ``generic flatness''.}
  Let \(S\) be a noetherian integral scheme, \(p : X \to S\) a finite-type morphism,
  and \(\F\) a coherent \(\OO_X\)-module. Then there exists a non-empty open
  subscheme \(U \subseteq S\) such that the restriction of \(\F\) to
  \(X_U = p^{-1}(U)\) is \(\OO_U\)-flat.
\end{theorem}

% SOURCE QUOTE PROOF: "The above theorem has the following consequence, which
% follows by restricting attention to a non-empty affine open subscheme of \(S\)."
% (Source: nitsure-hilbert-quot-src/nitsure-hilbert-quot.tex, L1776-L1777. The
% Nitsure proof is the single sentence quoted: deduce the geometric form from
% the algebraic generic flatness statement (\cref{thm:generic_flatness_algebraic})
% by restricting to a non-empty
% affine open of \(S\), where the situation reduces to the algebraic statement.)
\begin{proof}
  
  Restrict to a non-empty affine open \(\Spec A \subseteq S\) where \(A\) is a
  noetherian domain. The morphism \(p\) is of finite type, so locally on \(X\)
  above this affine \(\F\) is the sheaf \(\widetilde{M}\) for \(M\) a finite module
  over a finite-type \(A\)-algebra \(B\). \cref{thm:generic_flatness_algebraic}
  produces \(f \in A\), \(f \ne 0\), with \(M_f\) free over \(A_f\) on each such patch;
  shrinking to the basic open \(D(f) \subseteq \Spec A\) globalises by patching
  finitely many such opens (using noetherianity of \(X\) above the affine).
\end{proof}

\section{Existence of the flattening stratification}
\label{sec:flatstrat_exists}

\begin{theorem}

\leanok
[Existence of the flattening stratification]
  \label{thm:flattening_stratification_exists}
  \lean{AlgebraicGeometry.flatteningStratification}
  \uses{def:coherent_sheaf_flat, lem:flat_locus_open, lem:nonflat_locus_proper,
    lem:noetherian_induction_strata, thm:generic_flatness}

  % SOURCE: [Nitsure], \S4, Theorem ``flattening stratification''
  % (read from references/nitsure-hilbert-quot-src/nitsure-hilbert-quot.tex, L1811-L1841).
  % Parallel pointer: [Stacks Project], tag 052H (verbatim text not yet
  % retrieved into a local references/ file; see ``Notes for Plan Agent'').
  % SOURCE QUOTE:
  % "Let \(S\) be a noetherian scheme, and let
  %  \(\F\) be a coherent sheaf on the projective space \(\PP^n_S\)
  %  over \(S\). Then the set \(I\) of Hilbert polynomials of
  %  restrictions of \(\F\) to fibers of \(\PP^n_S\to S\)
  %  is a finite set. Moreover, for each \(f\in I\) there exist
  %  a locally closed subscheme \(S_f\) of \(S\),
  %  such that the following conditions are satisfied.
  %
  %  \textbf{(i) Point-set: } The underlying set \(|S_f|\) of \(S_f\) consists of
  %  all points \(s\in S\) where the Hilbert polynomial of the
  %  restriction of \(\F\) to \(\PP^n_s\) is \(f\). In particular,
  %  the subsets \(|S_f|\subset |S|\) are disjoint,
  %  and their set-theoretic union is \(|S|\).
  %
  %  \textbf{(ii) Universal property: }
  %  Let \(S'=\coprod S_f\) be the coproduct of the \(S_f\), and let
  %  \(i:S'\to S\) be the morphism induced by the inclusions
  %  \(S_f\hra S\). Then the sheaf \(i^*(\F)\) on \(\PP^n_{S'}\)
  %  is flat over \(S'\). Moreover, \(i:S'\to S\) has the universal property that
  %  for any morphism \(\varphi :T\to S\) the pullback \(\varphi^*(\F)\) on
  %  \(\PP^n_T\) is flat over \(T\) if and only if \(\varphi\) factors through
  %  \(i:S'\to S\).  The subscheme \(S_f\) is
  %  uniquely determined by the polynomial \(f\).
  %
  %  \textbf{(iii) Closure of strata: }
  %  Let the set \(I\) of Hilbert polynomials be given a total
  %  ordering, defined by putting \(f<g\) whenever \(f(n) < g(n)\) for all
  %  \(n \gg  0\). Then the closure in \(S\) of the subset \(|S_f|\)
  %  is contained in the union of all \(|S_g|\) where \(f\le g\)."
  \textit{Source: [Nitsure], \S4, Theorem ``flattening stratification'';
  cf.\ [Stacks Project], tag 052H.}
  Let \(S\) be a noetherian scheme and \(\F\) a coherent sheaf on the relative
  projective space \(\PP^n_S\) over \(S\).  Then:
  \begin{enumerate}[label=(\roman*)]
    \item The set \(I\) of Hilbert polynomials
      \(P_s(m) = \chi(\PP^n_s, \F_s(m))\) that arise as \(s\) ranges over \(S\) is
      finite.
    \item For each \(f \in I\) there exists a locally-closed subscheme
      \(S_f \hookrightarrow S\) whose underlying set \(|S_f|\) consists of exactly
      those \(s \in S\) with \(P_s = f\).  The collection \(\{|S_f|\}_{f \in I}\) is
      pairwise disjoint with set-theoretic union \(|S|\).  Forming the coproduct
      \(S' := \coprod_{f \in I} S_f\) and the morphism \(i : S' \to S\) assembled
      from the inclusions \(S_f \hookrightarrow S\), the pullback \(i^* \F\) on
      \(\PP^n_{S'}\) is flat over \(S'\), and \(i : S' \to S\) is universal with this
      property: for any morphism \(\varphi : T \to S\) such that \(\varphi^* \F\) on
      \(\PP^n_T\) is flat over \(T\), \(\varphi\) factors uniquely through \(i\).
      Each \(S_f\) is uniquely determined by \(f\).
    \item Order \(I\) by \(f < g \iff f(n) < g(n)\) for all \(n \gg 0\). Then the
      closure of \(|S_f|\) in \(|S|\) is contained in \(\bigcup_{g \ge f} |S_g|\).
  \end{enumerate}
\end{theorem}

% SOURCE QUOTE PROOF: "It is enough to prove the theorem for open subschemes of
% $S$ which cover $S$, as the resulting strata will then glue together by their
% universal property. [...] \textbf{General case: } We now allow the integer
% $n$ to be arbitrary. The idea of the proof is as follows: We show the
% existence of a stratification of $S$ which is a `g.c.d.' of the flattening
% stratifications for direct images $\pi_*\F(i)$ for all $i \ge N$ for some
% integer $N$ (where the flattening stratifications for $\pi_*\F(i)$ exist by
% case $n=0$ which we have treated above). This is the desired flattening
% stratification of $\F$ over $S$, as follows from Lemma \ref{graded module is flat}."
% (Source: nitsure-hilbert-quot-src/nitsure-hilbert-quot.tex, L1844-L1900; the
% full body L1844-L2090 is distributed across the sub-lemmas below, each of
% which carries its own verbatim % SOURCE QUOTE PROOF.)
\begin{proof}
  \uses{def:coherent_sheaf_flat, lem:flat_locus_open, lem:nonflat_locus_proper,
    lem:noetherian_induction_strata, thm:generic_flatness}
  Following [Nitsure], \S4, the construction reduces to the special case
  \(n = 0\), assembled over the relative twists \(\pi_* \F(N+i)\).  As the
  argument is local on \(S\) (the strata glue by their universal property, part
  (ii)), and as the three conclusions are proved through the sub-lemmas below,
  we assemble them in the order (i), (ii), (iii).

  \emph{Part (i): finiteness of the Hilbert-polynomial set.}
  By Lemma~\ref{lem:nonflat_locus_proper}, Noetherian induction on \(S\) (whose
  inductive step invokes generic flatness, Theorem~\ref{thm:generic_flatness},
  on the integral locally-closed strata cut out by the irreducible-component
  decomposition of \(S\)) produces finitely many reduced, pairwise disjoint,
  locally-closed subschemes \(V_i \hookrightarrow S\) set-theoretically covering
  \(|S|\) on each of which \(\F|_{\PP^n_{V_i}}\) is flat over \(\OO_{V_i}\) in
  the sense of Definition~\ref{def:coherent_sheaf_flat}.  On a flat family the
  Hilbert polynomial \(P_s(m) = \chi(\PP^n_s, \F_s(m))\) is locally constant, so
  it takes finitely many values on each quasi-compact \(V_i\); the finite union
  over the finitely many \(V_i\) shows that the set \(I\) of Hilbert polynomials
  occurring on \(S\) is finite.

  \emph{Part (ii): existence of the strata and their universal property.}
  Working over the flat pieces \(V_i\) of part (i), semicontinuity
  (cohomology-and-base-change) yields a uniform integer \(N\) such that, for all
  \(s \in S\) and \(m \ge N\), the higher cohomologies \(H^r(\PP^n_s, \F_s(m))\)
  vanish for \(r \ge 1\) and the base-change homomorphism
  \((\pi_* \F(m))|_s \to H^0(\PP^n_s, \F_s(m))\) is an isomorphism.  With this
  \(N\) fixed, Lemma~\ref{lem:noetherian_induction_strata} applies the \(n = 0\)
  stratification mechanism (Lemma~\ref{lem:flat_locus_open}) successively to the
  coherent sheaves \(E_i = \pi_* \F(N+i)\) on \(S\) for \(i = 0, \dots, n\),
  refining stratum by stratum.  After \(n+1\) steps this produces a finite
  family of locally-closed subschemes \(W_{e_0, \dots, e_n} \subseteq S\),
  universal for the property that each \(E_i\) pulls back to a locally-free sheaf
  of constant rank \(e_i\); re-indexing the tuple \((e_0, \dots, e_n)\) by the
  unique numerical polynomial \(f\) of degree \(\le n\) with \(f(N+i) = e_i\)
  gives the family \(\{W_f\}_{f \in I}\).  At any \(s \in W_f\) the base-change
  isomorphism and the cohomology vanishing force \(P_s = f\) (since \(P_s\) has
  degree \(\le n\) and is determined by its values at \(N, \dots, N+n\)), so
  \(|W_f| = \{s : P_s = f\}\); these subsets are disjoint with union \(|S|\).
  The required locally-closed subscheme \(S_f \subseteq W_f\) is the closed
  subscheme cut out by the further local-freeness conditions
  ``\(\pi_* \F(N+i)|_{W_f}\) is locally free of rank \(f(N+i)\) for all
  \(i \ge 0\)'' (a chain of coherent ideal sheaves that stabilises by
  Noetherianity), and \(\F|_{\PP^n_{S_f}}\) is flat over \(S_f\)
  (Definition~\ref{def:coherent_sheaf_flat}) because the graded module
  \(\bigoplus_{i \gg 0} \pi_* \F(i)|_{S_f}\) is locally free.  Forming
  \(S' = \coprod_{f \in I} S_f\) and the assembled morphism \(i : S' \to S\),
  the pullback \(i^* \F\) is \(S'\)-flat; and any \(\varphi : T \to S\) with
  \(\varphi^* \F\) flat over \(T\) has locally constant Hilbert polynomial, hence
  decomposes \(T = \coprod_f T_f\) into the open-closed loci where the polynomial
  is \(f\), each factoring uniquely through the corresponding \(S_f\) by the
  local-freeness universal property of Lemma~\ref{lem:flat_locus_open}.  This
  gives the unique factorisation through \(i\), and shows each \(S_f\) is
  uniquely determined by \(f\).

  \emph{Part (iii): closure of strata.}
  Order \(I\) by \(f < g \iff f(n) < g(n)\) for all \(n \gg 0\).  Since each
  \(S_f\) is a stratum of the \(n = 0\) flattening stratification applied to a
  single direct image \(\pi_* \F(p)\) for \(p \gg 0\) separating all polynomials
  in \(I\), the closure-of-strata assertion reduces to the upper-semicontinuity
  statement already established in the \(n = 0\) case
  (Lemma~\ref{lem:flat_locus_open}): the closure of \(|S_f|\) in \(|S|\) is
  contained in \(\bigcup_{g \ge f} |S_g|\).
\end{proof}

The proof is long; we split it into sub-lemmas below.  The first sub-lemma is
the special case \(n = 0\) (when \(\PP^n_S = S\) and \(\F\) is just a coherent sheaf
on \(S\)), which already contains the flat-locus-as-locally-closed-subscheme
mechanism. The remaining inputs are (A) the Noetherian-induction reduction to
the flat case via generic flatness; (B) uniform vanishing of higher cohomology
of \(\F(m)\) for \(m \gg 0\) on \(\F_{V_i}\) for the flat pieces; and (C) base-change
isomorphisms for the direct images \(\pi_* \F(m)\), which allow comparing
strata for different twists.

\begin{lemma}
[Flat locus on \(S\) for a coherent sheaf is a locally-closed stratification]
  \label{lem:flat_locus_open}
  \lean{AlgebraicGeometry.TODO.flatLocusOpen}

  % SOURCE: [Nitsure], \S4, Theorem ``flattening stratification'', Special case
  % \(n=0\) (read from references/nitsure-hilbert-quot-src/nitsure-hilbert-quot.tex,
  % L1849-L1885).
  % SOURCE QUOTE:
  % "\textbf{Special case: } Let \(n=0\), so that \(\PP^n_S=S\).
  %  For any \(s\in S\), the \textbf{fiber} \(\F|_s\) of \(\F\) over \(s\)
  %  will mean the pull-back of \(\F\) to the subscheme
  %  \(\spec \kappa(s)\) [...]. The Hilbert polynomial of the
  %  restriction of \(\F\) to the fiber over \(s\) is the degree \(0\)
  %  polynomial \(e \in \Q[\lambda]\), where \(e = \dim_{\kappa(s)}\F|_s\).
  %
  %  By Nakayama lemma, any basis of \(\F|_s\) prolongs to a neighbourhood \(U\) of \(s\)
  %  to give a set of generators for \(F|_U\). Repeating this argument,
  %  we see that there exists a smaller neighbourhood \(V\) of \(s\) in which
  %  there is a right-exact sequence
  %  ${\mathcal O}_V^{\oplus m} \stackrel{\psi}{\to} {\mathcal O}_V^{\oplus e}
  %  \stackrel{\phi}{\to} \F \to 0$.
  %  Let \(I_e\subset {\mathcal O}_V\) be the ideal sheaf
  %  formed by the entries of the \(e\times m\) matrix
  %  \((\psi_{i,j})\) [...].
  %  Let \(V_e\) be the closed subscheme of \(V\) defined by \(I_e\).
  %  For any morphism of schemes
  %  \(f: T\to V\), the pull-back sequence [...]
  %  is exact, by right-exactness of tensor products.
  %  Hence the pull-back \(f^*\F\) is a locally free \({\mathcal O}_T\)-module of
  %  rank \(e\) if and only if \(f^*\psi =0\), that is,
  %  \(f\) factors via the subscheme \(V_e\hra V\) defined by the vanishing of all
  %  entries \(\psi_{i,j}\).
  %  Thus we have proved assertions \textbf{(i)} and \textbf{(ii)} of the theorem.
  %
  %  As the rank of the matrix \((\psi_{i,j})\) is lower semi-continuous,
  %  it follows that the function \(e\) is upper semi-continuous, which
  %  proves the assertion \textbf{(iii)} of the theorem, completing its proof
  %  when \(n=0\)."
  \textit{Source: [Nitsure], \S4 (special case \(n=0\) of the main theorem).}
  Let \(S\) be a noetherian scheme and let \(\F\) be a coherent \(\OO_S\)-module.
  For each integer \(e \ge 0\), there exists a locally-closed subscheme
  \(S_e \hookrightarrow S\) whose underlying set is
  \(\{s \in S : \dim_{\kappa(s)} \F|_s = e\}\). The family \(\{S_e\}_e\) is a
  finite, mutually disjoint, set-theoretically covering stratification of \(S\),
  and a morphism \(f : T \to S\) pulls back \(\F\) to a locally-free
  \(\OO_T\)-module of rank \(e\) if and only if \(f\) factors through \(S_e\).  The
  function \(s \mapsto \dim_{\kappa(s)} \F|_s\) is upper-semicontinuous on \(S\).
\end{lemma}

\begin{proof}
  
  By [Nitsure], \S4 (verbatim above): for \(s \in S\), Nakayama produces a
  neighbourhood \(V\) of \(s\) and a right-exact local presentation
  \(\OO_V^{\oplus m} \xrightarrow{\psi} \OO_V^{\oplus e} \to \F|_V \to 0\) where
  \(e = \dim_{\kappa(s)} \F|_s\).  The closed subscheme \(V_e \hookrightarrow V\)
  defined by the ideal sheaf generated by the matrix entries of \(\psi\)
  represents the locus where \(\F\) becomes locally free of rank \(e\): by
  right-exactness of tensor product, a base change \(f : T \to V\) has
  \(f^* \F\) locally free of rank \(e\) iff \(f^* \psi = 0\), iff \(f\) factors
  through \(V_e\).  Glueing the \(V_e\) over the cover of \(S\) by such
  neighbourhoods (and intersecting with \(V \setminus \bigcup_{e' > e} V_{e'}\)
  to get the strictly-equal-to-\(e\) stratum) produces the locally-closed
  \(S_e\). Upper-semicontinuity of \(e\) follows because the rank of \(\psi\) is
  lower-semicontinuous.
\end{proof}

\begin{lemma}
[Noetherian induction reduction to the flat case]
  \label{lem:nonflat_locus_proper}
  \lean{AlgebraicGeometry.TODO.nonflatLocusProper}
  \uses{thm:generic_flatness}

  % SOURCE: [Nitsure], \S4, Theorem ``flattening stratification'', General case
  % opening paragraphs (read from
  % references/nitsure-hilbert-quot-src/nitsure-hilbert-quot.tex, L1888-L1929).
  % SOURCE QUOTE:
  % "\textbf{General case: } We now allow the integer \(n\) to be arbitrary.
  %  The idea of the proof is as follows:
  %  We show the existence of a stratification
  %  of \(S\) which is a `g.c.d.' of the flattening
  %  stratifications for direct images \(\pi_*\F(i)\) for all \(i \ge N\) for some
  %  integer \(N\) [...].
  %
  %  As \(S\) is noetherian, it is a finite union of irreducible components,
  %  and these are closed in \(S\).
  %  Let \(Y\) be an irreducible component of \(S\), and let \(U\) be the
  %  non-empty open subset of \(Y\) which consists of all points which do not
  %  lie on any other irreducible component of \(S\).
  %  Let \(U\) be given the reduced subscheme
  %  structure. Note that this makes \(U\) an integral scheme, which is a
  %  locally closed subscheme of \(S\).
  %  By Theorem \ref{generic flatness} on generic flatness,
  %  \(U\) has a non-empty open subscheme \(V\) such that
  %  the restriction of \(\F\) to \(\PP^n_V\) is flat over \({\mathcal O}_V\).
  %  Now repeating the argument with \(S\) replaced by its reduced
  %  closed subscheme \(S-V\),
  %  it follows by noetherian induction on \(S\) that there exist
  %  finitely many reduced, locally closed, mutually disjoint subschemes
  %  \(V_i\) of \(S\) such that set-theoretically
  %  \(|S|\) is the union of the \(|V_i|\) and
  %  the restriction of \(\F\) to \(\PP^n_{V_i}\) is flat over \({\mathcal O}_{V_i}\).
  %  As each \(V_i\) is a noetherian scheme, and as the Hilbert polynomials
  %  are locally constant for a flat family of sheaves, it follows
  %  that only finitely many polynomials occur in \(V_i\)
  %  in the family of Hilbert polynomials \(P_s(m) = \chi(\PP^n_s,\F_s(m))\)
  %  as \(s\) varies over points of \(V_i\). This allows us to conclude
  %  the following:
  %
  %  \textbf{(A) } Only finitely many distinct Hilbert polynomials
  %  \(P_s(m) = \chi(\PP^n_s,\F_s(m))\) occur, as \(s\)
  %  varies over all of \(S\)."
  \textit{Source: [Nitsure], \S4 (general case, reduction by Noetherian induction).}
  Let \(S\) be a noetherian scheme, \(\pi : \PP^n_S \to S\) the relative projective
  space, and \(\F\) a coherent sheaf on \(\PP^n_S\).  Then there exist finitely many
  reduced, locally-closed, pairwise disjoint subschemes \(V_i \hookrightarrow S\)
  whose set-theoretic union is \(|S|\) and such that \(\F|_{\PP^n_{V_i}}\) is
  flat over \(\OO_{V_i}\) for each \(i\). Consequently, only finitely many
  Hilbert polynomials \(P_s(m) = \chi(\PP^n_s, \F_s(m))\) occur as \(s\) varies
  over \(S\).
\end{lemma}

\begin{proof}
  
  By [Nitsure], \S4: \(S\) being noetherian, decompose \(|S| = \bigcup Y_j\) into
  finitely many irreducible components, each closed.  For an irreducible
  component \(Y\), let \(U \subseteq Y\) be the open subset not lying on any other
  component, given the reduced subscheme structure (an integral, locally-closed
  subscheme of \(S\)).  By \cref{thm:generic_flatness} applied to
  \(\F|_{\PP^n_U}\), there is a non-empty open \(V \subseteq U\) over which \(\F\) is
  flat.  Replace \(S\) by the reduced closed complement \(S \setminus V\) and
  repeat. Noetherianity of \(S\) terminates the induction in finitely many
  steps, producing the \(V_i\).  On each \(V_i\) the Hilbert polynomial is locally
  constant (flat families have locally constant Hilbert polynomial), hence
  takes finitely many values on the noetherian (hence quasi-compact) scheme
  \(V_i\); the finite union of finite sets is finite.
\end{proof}

\begin{lemma}
[Assembly of the strata by Noetherian induction]
  \label{lem:noetherian_induction_strata}
  \lean{AlgebraicGeometry.TODO.noetherianInductionStrata}
  \uses{lem:flat_locus_open, lem:nonflat_locus_proper, thm:flat_base_change_cohomology}

  % SOURCE: [Nitsure], \S4, Theorem ``flattening stratification'', Assembly via
  % direct images \(E_i = \pi_*\F(N+i)\) (read from
  % references/nitsure-hilbert-quot-src/nitsure-hilbert-quot.tex, L1992-L2090).
  % SOURCE QUOTE:
  % "Let \(\pi : \PP^n_S \to S\) denote the projection. Consider the
  %  coherent sheaves \(E_0, \ldots, E_n\) on \(S\), defined by
  %  \(E_i = \pi_*\F(N+i)\mbox{ for } i=0,\ldots,n.\)
  %  By applying the special case of of the theorem (where the relative dimension
  %  \(n\) of \(\PP^n_S\) is \(0\)) to the sheaf \(E_0\) on \(\PP^0_S=S\), we get
  %  a stratification \((W_{e_0})\) of \(S\)
  %  indexed by integers \(e_0\), such that for any morphism
  %  \(f: T\to S\) the pull-back \(f^*E_0\) is a locally free \({\mathcal O}_T\)-module
  %  of rank \(e_0\) if and only if \(f\) factors via \(W_{e_0} \hra S\).
  %  Next, for each stratum \(W_{e_0}\), we take the
  %  flattening stratification \((W_{e_0,e_1})\) for \(E_1|_{W_{e_0}}\), and so on.
  %  Thus in \(n+1\) steps, we obtain finitely many
  %  locally closed subschemes
  %  \(W_{e_0,\ldots,e_n} \subset S\)
  %  such that for any morphism \(f:T\to S\) the pull-back
  %  \(f^*E_i\) for \(i=0,\ldots,n\) is a locally free \({\mathcal O}_T\)-modules of
  %  of constant rank \(e_i\) if and only if \(f\) factors
  %  via \(W_{e_0,\ldots,e_n} \hra S\). [...]
  %  This shows that at any point \(s\in W_f\), the Hilbert polynomial
  %  \(P_s(m)\) equals \(f\). The desired locally closed subscheme \(S_f \subset S\),
  %  whose existence is asserted by the theorem,
  %  will turn out to be a certain closed subscheme \(S_f \subset W_f\) whose
  %  underlying subset is all of \(|W_f|\). [...]
  %  As \(S_f\) is the flattening stratification for \(\pi_*\F(p)\),
  %  the property (iii) of the theorem follows from the corresponding property in
  %  the case \(n=0\), applied to the sheaf \(\pi_*\F(p)\) on \(S\)."
  \textit{Source: [Nitsure], \S4 (assembly step combining the \(n=0\) special case
  applied to direct images \(\pi_*\F(N+i)\)).}
  Let \(S\), \(\pi : \PP^n_S \to S\), \(\F\) be as in
  \cref{thm:flattening_stratification_exists}.  There exists an integer
  \(N \ge 0\) such that for every \(s \in S\) the higher cohomologies
  \(H^r(\PP^n_s, \F_s(m))\) vanish for all \(r \ge 1\), \(m \ge N\), and the base-change
  homomorphism \((\pi_* \F(m))|_s \to H^0(\PP^n_s, \F_s(m))\) is an isomorphism
  for \(m \ge N\), \(s \in S\).  Iterating the \(n=0\) stratification
  (\cref{lem:flat_locus_open}) on the coherent sheaves \(E_i := \pi_* \F(N+i)\)
  for \(i = 0, \dots, n\) produces a finite collection of locally-closed
  subschemes \(W_{e_0, \dots, e_n} \subseteq S\) indexed by tuples
  \((e_0, \dots, e_n) \in \Z^{n+1}\). Re-indexing each tuple by the unique
  numerical polynomial \(f \in \Q[\lambda]\) of degree \(\le n\) with
  \(f(N+i) = e_i\) gives a finite stratification \(\{W_f\}\) of \(S\), and the
  locally-closed subschemes \(S_f \subseteq W_f\) of
  \cref{thm:flattening_stratification_exists} are obtained by further
  intersecting with the closed-subscheme conditions that
  \(\pi_* \F(N+i)|_{W_f}\) be locally free of rank \(f(N+i)\) for all \(i \ge 0\)
  (terminating by Noetherianity).
\end{lemma}

% SOURCE QUOTE PROOF: full proof is the long body of [Nitsure], \S4, Theorem
% ``flattening stratification'', spanning L1844-L2090 (~250 lines: special
% case, Noetherian induction, statements (A)/(B)/(C), then the assembly via
% direct images and the closed-subscheme construction inside each \(W_f\)).
% Verbatim retrieved per-sub-lemma above; the structural restatement below in
% project notation is the canonical Nitsure-style restated proof.
\begin{proof}
  
  By \cref{lem:nonflat_locus_proper}, the family of Hilbert polynomials
  \(\{P_s\}\) on \(S\) is finite, and there is a finite stratification
  \(\{V_i\}\) of \(S\) on each of which \(\F\) is flat with locally-constant
  Hilbert polynomial.  Apply semi-continuity (cohomology and base change,
  [Hartshorne], III.12 / [Stacks Project], tag 02KH for the \(H^0\)-version) on
  the noetherian \(V_i\) to obtain a uniform integer \(N_1\) with
  \(R^r \pi_* \F(m) = 0\) for \(r \ge 1\), \(m \ge N_1\), and \(H^r(\PP^n_s, \F_s(m)) = 0\)
  for all \(s \in S\), \(r \ge 1\), \(m \ge N_1\) (statement (B) of Nitsure's proof).
  Choose \(r_i \ge N_1\) on each \(V_i\) such that the base-change comparison
  \((\pi_* \F(m))|_{V_i} \to (\pi_i)_* \F_{V_i}(m)\) is an isomorphism for
  \(m \ge r_i\); composing with the cohomology-and-base-change isomorphism
  \(((\pi_i)_* \F_{V_i}(m))|_s \xrightarrow{\sim} H^0(\PP^n_s, \F_s(m))\) for
  the flat family on \(V_i\) yields uniform \(N = \max r_i\) such that
  \((\pi_* \F(m))|_s \to H^0(\PP^n_s, \F_s(m))\) is an isomorphism for \(m \ge N\)
  and all \(s \in S\) (statement (C)).

  Apply \cref{lem:flat_locus_open} sequentially to the coherent sheaves
  \(E_i = \pi_* \F(N+i)\) on \(S\) for \(i = 0, \dots, n\): \(E_0\) produces a
  stratification \(\{W_{e_0}\}\) of \(S\); on each \(W_{e_0}\) apply the lemma to
  \(E_1|_{W_{e_0}}\) to refine to \(\{W_{e_0, e_1}\}\), and continue. After \(n+1\)
  steps the \(W_{e_0, \dots, e_n}\) form a finite locally-closed stratification
  of \(S\) universal for the property ``each \(E_i\) pulls back to a locally-free
  \(\OO_T\)-module of constant rank \(e_i\)''.  Re-index \((e_0, \dots, e_n)\) by the
  unique numerical polynomial \(f\) of degree \(\le n\) with \(f(N+i) = e_i\) to get
  the family \(\{W_f\}\) indexed by polynomials.  At any \(s \in W_f\), statement
  (C) and vanishing (B) force \(P_s = f\) (since \(P_s\) has degree \(\le n\) and is
  determined by its values at \(N, N+1, \dots, N+n\)).

  The Hilbert-polynomial-coincidence locus \(|W_f|\) may have non-flat scheme
  structure; refine to \(S_f \subseteq W_f\) by intersecting with the
  closed-subscheme conditions ``\(\pi_* \F(N+i)|_{W_f}\) is locally-free of rank
  \(f(N+i)\)'' for all \(i \ge 0\).  These conditions are cut out by an increasing
  chain of coherent ideal sheaves \(I_0 \subseteq I_0 + I_1 \subseteq \cdots\) on
  \(W_f\); Noetherianity terminates the chain in finitely many steps, defining a
  coherent ideal sheaf \(I\) and a closed subscheme \(S_f \subseteq W_f\) with
  \(|S_f| = |W_f|\) on which \(\pi_* \F(N+i)\) is locally-free of rank \(f(N+i)\)
  for all \(i \ge 0\). The base-change-without-flatness lemma (cf.\ [Nitsure],
  \S3) yields some \(N' \ge N\) such that for \(i \ge N'\),
  \((\pi_* \F(i))|_{S_f} \to (\pi_{S_f})_* \F_{S_f}(i)\) is an isomorphism, hence
  the graded \(A\)-module \(\bigoplus_{i \ge N'} (\pi_{S_f})_* \F_{S_f}(i)\) is
  locally-free on \(S_f\), hence \(\F_{S_f}\) is flat over \(S_f\). The
  closure-of-strata property (iii) reduces to the corresponding property in
  the \(n=0\) case for \(\pi_* \F(p)\) on \(S\) at any \(p \ge N\) separating all
  Hilbert polynomials.
\end{proof}

\subsection{Lean-side encodings of the sub-lemmas}
\label{sec:flatstrat_lean_helpers}

The three sub-lemmas above are realised in Lean by the following helper
declarations, which carry the substantive type signatures consumed by the
existence theorem. Each is an internal project helper with no external source
beyond Nitsure~\S4 already cited above.

\begin{lemma}

\leanok
[Flat-locus stratification, special case \(n = 0\)]
  \label{lem:flat_locus_stratification_lean}
  \lean{AlgebraicGeometry.flatLocusStratification}
  \uses{def:coherent_sheaf_flat}
  Let \(S\) be a locally noetherian scheme and \(\F\) a coherent \(\OO_S\)-module.
  Then there is a family of immersions \(\iota_e : S_e \to S\) indexed by
  \(e \in \mathbb{N}\) whose images are pairwise disjoint and set-theoretically
  cover \(|S|\), and such that for each \(e\) the pullback of \(\F\) along
  \(\iota_e\) is flat over \(S_e\) in the sense of
  \cref{def:coherent_sheaf_flat}.  This is the Lean encoding of the special
  case \(n = 0\) of the flattening stratification
  (\cref{lem:flat_locus_open}), the rank index \(e\) playing the role of
  \(\dim_{\kappa(s)} \F|_s\).
\end{lemma}

\begin{proof}
  Proved directly in Lean.
\end{proof}

\begin{lemma}

\leanok
[Noetherian-induction reduction to the flat case]
  \label{lem:flat_locus_reduction_lean}
  \lean{AlgebraicGeometry.flatLocusReduction}
  \uses{def:coherent_sheaf_flat, thm:generic_flatness}
  Let \(S\) be a locally noetherian scheme, \(\pi : X \to S\) a proper morphism,
  and \(\F\) a coherent \(\OO_X\)-module. Then there exist finitely many
  immersions \(\iota_i : V_i \to S\) with pairwise disjoint images
  set-theoretically covering \(|S|\), such that for each \(i\) the pullback of
  \(\F\) to \(X \times_S V_i\) is flat over \(V_i\) in the sense of
  \cref{def:coherent_sheaf_flat}.  This is the Lean encoding of the
  Noetherian-induction reduction (\cref{lem:nonflat_locus_proper}), obtained by
  peeling off non-empty flat open patches via generic flatness
  (\cref{thm:generic_flatness}).
\end{lemma}

\begin{proof}
  Proved directly in Lean.
\end{proof}

\begin{lemma}

\leanok
[Assembly of the Hilbert-polynomial-indexed strata]
  \label{lem:flat_locus_assembly_lean}
  \lean{AlgebraicGeometry.flatLocusAssembly}
  \uses{def:coherent_sheaf_flat, lem:flat_locus_stratification_lean,
    lem:flat_locus_reduction_lean}
  Let \(S\) be a locally noetherian scheme, \(\pi : X \to S\) a proper morphism,
  and \(\F\) a coherent \(\OO_X\)-module. Then there exist a finite index set
  \(I\), immersions \(\iota_f : S_f \to S\) (\(f \in I\)) with pairwise disjoint
  images set-theoretically covering \(|S|\), and a label map
  \(P : I \to (\mathbb{N} \to \Z)\) with \(f = g \iff P_f = P_g\), such that for
  each \(f\) the pullback of \(\F\) to \(X \times_S S_f\) is flat over \(S_f\) in
  the sense of \cref{def:coherent_sheaf_flat}.  This is the Lean encoding of the
  assembly step (\cref{lem:noetherian_induction_strata}): the special-case
  stratification (\cref{lem:flat_locus_stratification_lean}) is iterated on the
  direct images \(\pi_* \F(N+i)\) to refine the reduction
  (\cref{lem:flat_locus_reduction_lean}) into the Hilbert-polynomial-indexed
  family, with \(P\) recording the numerical-polynomial label of each stratum.
\end{lemma}

\begin{proof}
  Proved directly in Lean.
\end{proof}

\section{Universal property of the strata}
\label{sec:flatstrat_universal}

\begin{theorem}

\leanok
[Universal property of the flattening stratification]
  \label{thm:flattening_stratification_universal}
  \lean{AlgebraicGeometry.flatteningStratification_universal}
  \uses{thm:flattening_stratification_exists, lem:noetherian_induction_strata}

  % SOURCE: [Nitsure], \S4, Theorem ``flattening stratification'', part (ii)
  % universal property
  % (read from references/nitsure-hilbert-quot-src/nitsure-hilbert-quot.tex, L1826-L1834).
  % SOURCE QUOTE:
  % "\textbf{(ii) Universal property: }
  %  Let \(S'=\coprod S_f\) be the coproduct of the \(S_f\), and let
  %  \(i:S'\to S\) be the morphism induced by the inclusions
  %  \(S_f\hra S\). Then the sheaf \(i^*(\F)\) on \(\PP^n_{S'}\)
  %  is flat over \(S'\). Moreover, \(i:S'\to S\) has the universal property that
  %  for any morphism \(\varphi :T\to S\) the pullback \(\varphi^*(\F)\) on
  %  \(\PP^n_T\) is flat over \(T\) if and only if \(\varphi\) factors through
  %  \(i:S'\to S\).  The subscheme \(S_f\) is
  %  uniquely determined by the polynomial \(f\)."
  \textit{Source: [Nitsure], \S4 (part (ii) of the main theorem); cf.\ [Stacks
  Project], tag 052H.}
  Let \(S\), \(\F\), and \(\{S_f\}_{f \in I}\) be as in
  \cref{thm:flattening_stratification_exists}.  Set
  \(S' := \coprod_{f \in I} S_f\) with the coproduct morphism
  \(i : S' \to S\) assembled from the locally-closed inclusions.  Then \(i^* \F\)
  is flat over \(S'\), and the pair \((S', i)\) is universal: for any morphism
  \(\varphi : T \to S\) of noetherian schemes (one may take \(T\) arbitrary, the
  Noetherian hypothesis on \(S\) being preserved), the pullback
  \(\varphi^* \F\) on \(\PP^n_T\) is \(T\)-flat if and only if \(\varphi\) factors
  uniquely through \(i\). In particular each \(S_f\) is uniquely determined by \(f\).
\end{theorem}

% SOURCE QUOTE PROOF: "We now show that the morphism \(\coprod_f S_f \,\to \,S\)
% indeed has the property (ii) of the theorem.
% By Lemma \ref{base change without flatness},
% there exists some \(N' \ge N\) such that for all \(i\ge N'\),
% the base-change
% \((\pi_*\F(i))|_{S_f} \to (\pi_{S_f})_*\F_{S_f}(i)\) is an isomorphism for
% each \(S_f\). Therefore
% \(\F_{S_f}\) is flat over \(S_f\) by Lemma \ref{graded module is flat},
% as the direct images \(\pi_*\F(i)\) for all \(i \ge N'\) are locally
% free over \(S_f\). Conversely, if \(\phi : T\to S\) is a morphism such
% that \(\F_T\) is flat, then the Hilbert polynomial is locally
% constant over \(T\). Let \(T_f\) be the open and closed subscheme of \(T\)
% where the Hilbert polynomial is \(f\). Clearly, the set
% map \(|T_f| \to |S|\) factors via \(|S_f|\). But as the direct images
% \({\pi_T}_*\F_T(i)\) are locally free of rank \(f(i)\) on \(T_f\), it
% follows in fact that the schematic morphism \(T_f \to S\) factors
% via \(S_f\), proving the property (ii) of the theorem."
% (Source: nitsure-hilbert-quot-src/nitsure-hilbert-quot.tex, L2065-L2082.)
\begin{proof}
  
  Flatness of \(i^* \F\) over \(S'\) was established at the end of the proof of
  \cref{thm:flattening_stratification_exists} (the
  base-change-without-flatness lemma plus the graded-module-flatness criterion
  applied to the coherent sheaves \(\pi_* \F(N+i)\) on each \(S_f\)).
  Conversely, if \(\varphi : T \to S\) satisfies \(\varphi^* \F\) on \(\PP^n_T\) is
  \(T\)-flat, then \(T\)-flatness of the family forces the Hilbert polynomial to be
  locally constant on \(T\).  Decompose \(T = \coprod_f T_f\) into the open-closed
  subschemes where the Hilbert polynomial is \(f\).  The set map \(|T_f| \to |S|\)
  factors through \(|S_f|\); further, since the direct images
  \((\pi_T)_* \varphi^* \F(i)\) are locally-free of rank \(f(i)\) on \(T_f\) for
  \(i \ge N'\), the universal property of \(S_f \subseteq W_f\) as the closed
  subscheme cut out by these local-freeness conditions
  (\cref{lem:flat_locus_open}) yields a unique scheme
  morphism \(T_f \to S_f\) extending the set map. Assembling over \(f\) gives the
  unique factorisation \(\varphi = i \circ \tilde \varphi\).
\end{proof}

\section{Application: relative curves}
\label{sec:flatstrat_curve_application}

In the Route~A consumer setting, the input to the flattening-stratification
machinery is the pullback \(C \times_k T \to T\) where \(C\) is a smooth proper
curve over a field \(k\) and \(T\) is any \(k\)-scheme.  The hypotheses of
Nitsure's \cref{thm:flattening_stratification_exists} demand
\emph{projective} relative dimension, not merely \emph{proper}; we record that
no extra work is required to instantiate the theorem in our setting.

\begin{lemma}
[Smooth proper curve over \(k\) is projective]
  \label{lem:smooth_proper_curve_projective}
  \lean{AlgebraicGeometry.TODO.smoothProperCurveProjective}
  % SOURCE: [Hartshorne], II.6.7 / II.7.6 (curve \(\Rightarrow\) quasi-projective;
  % proper \(\Rightarrow\) projective for curves), and standard consequences in
  % chapter III. The precise statement ``every smooth proper curve over a field
  % is projective'' is one of the classical equivalences for \(\dim 1\) and is
  % textbook material; this block is project-bespoke (no verbatim source
  % quote is reproduced here as the result is folklore-level and cited only
  % structurally for the Route~A reduction).
  \textit{Source: classical; cf.\ [Hartshorne], II.6.7 \& II.7.6 (curves).}
  Let \(C\) be a smooth proper curve over a field \(k\).  Then \(C \to \Spec k\) is
  projective.  Consequently, for any noetherian \(k\)-scheme \(T\), the relative
  curve \(C \times_k T \to T\) is projective in the sense of Grothendieck (there
  exists a relatively very ample line bundle on \(C \times_k T\) over \(T\)).
\end{lemma}

\begin{proof}
  Projectivity of a smooth proper curve over a field is classical
  ([Hartshorne], II.6.7 / II.7.6; alternatively: every divisor of large
  positive degree on a smooth proper curve \(C\) is very ample, providing the
  embedding \(C \hookrightarrow \PP^N_k\) for some \(N\)). Pulling the ample line
  bundle back along \(C \times_k T \to C\) produces a relatively very ample line
  bundle on \(C \times_k T\) over \(T\).
\end{proof}

\begin{corollary}

\leanok
[Flattening stratification for a coherent sheaf on a relative curve]
  \label{cor:flattening_stratification_curve}
  \lean{AlgebraicGeometry.flatteningStratification.ofCurve}
  \uses{lem:smooth_proper_curve_projective, thm:flattening_stratification_exists}
  Let \(k\) be a field, \(C\) a smooth proper curve over \(k\), and \(T\) a noetherian
  \(k\)-scheme.  Let \(\F\) be a coherent sheaf on \(C \times_k T\). Then the
  flattening-stratification conclusion of
  \cref{thm:flattening_stratification_exists} applies: there is a finite
  set \(I\) of Hilbert polynomials and locally-closed subschemes
  \(T_f \hookrightarrow T\) for \(f \in I\), set-theoretically covering \(T\)
  disjointly, such that (a) \(\F|_{C \times_k T_f}\) is flat over \(T_f\) for
  each \(f\), and (b) any morphism \(\varphi : T' \to T\) such that
  \(\varphi^* \F\) on \(C \times_k T'\) is \(T'\)-flat factors uniquely through some
  \(T_f\).
\end{corollary}

\begin{proof}
  
  By \cref{lem:smooth_proper_curve_projective}, \(C \times_k T \to T\) is
  projective, hence factors through a closed immersion into \(\PP^n_T\) for some
  \(n \ge 0\).  Push \(\F\) forward along the closed immersion (extending by zero
  outside the image) to a coherent sheaf \(\tilde \F\) on \(\PP^n_T\); this
  operation preserves \(T\)-flatness in both directions because the closed
  immersion is a finite morphism.  Apply
  \cref{thm:flattening_stratification_exists} to \(\tilde \F\) on
  \(\PP^n_T\) to obtain \(\{T_f\}\) with the stated universal property.
\end{proof}

\section{Mathlib status}
\label{sec:flatstrat_mathlib_status}

Mathlib (master \texttt{b80f227}) provides:
\begin{itemize}
  \item \texttt{Module.Flat} (\texttt{Mathlib.RingTheory.Flat.Basic}) --- the
    module-theoretic notion of flatness, applicable affine-locally to the
    stalks \(\F_x\).
  \item \texttt{Mathlib.AlgebraicGeometry.Morphisms.Flat} ---
    \texttt{IsFlat} as a morphism property and basic stability lemmas; this is
    flatness of the structure morphism, \emph{not} flatness of a quasi-coherent
    module over the base.
  \item \texttt{Module.FaithfullyFlat} and the descent infrastructure that
    consumes it.
\end{itemize}

\noindent What Mathlib does \emph{not} yet have, and that this chapter will
introduce in the new file
\texttt{AlgebraicJacobian/Picard/FlatteningStratification.lean}:
\begin{itemize}
  \item A predicate \texttt{AlgebraicGeometry.Scheme.CoherentSheafFlat}
    capturing \(S\)-flatness of a coherent \(\OO_X\)-module on a finite-type
    \(S\)-scheme (\cref{def:coherent_sheaf_flat}). This is the
    sheaf-level lift of \texttt{Module.Flat} via the affine-local
    characterization in the second paragraph of the definition.
  \item The constructive existence proof
    \texttt{AlgebraicGeometry.flatteningStratification}
    (\cref{thm:flattening_stratification_exists}) verbatim from
    Nitsure~\S4.  The proof decomposes into the three sub-lemmas
    \texttt{flatLocusStratification} (\cref{lem:flat_locus_stratification_lean}),
    \texttt{flatLocusReduction} (\cref{lem:flat_locus_reduction_lean}, the
    Noetherian-induction reduction via generic flatness), and
    \texttt{flatLocusAssembly} (\cref{lem:flat_locus_assembly_lean}).
  \item The universal-property packaging
    \texttt{AlgebraicGeometry.flatteningStratification\_universal}
    (\cref{thm:flattening_stratification_universal}) as a
    \texttt{CategoryTheory.Functor.RepresentableBy}-style universal arrow.
  \item The relative-curve specialisation
    \texttt{AlgebraicGeometry.flatteningStratification.ofCurve}
    (\cref{cor:flattening_stratification_curve}) for the Route~A
    consumer.
\end{itemize}

Inputs from Mathlib that the proof consumes:
\begin{itemize}
  \item \texttt{CategoryTheory.Limits.Cofan} / \texttt{Sigma} for the coproduct
    \(\coprod_f S_f\).
  \item \texttt{AlgebraicGeometry.IsLocallyClosedImmersion} (and the locally-closed
    subscheme API) for the strata \(S_f \hookrightarrow S\).
  \item \texttt{IsNoetherian} / \texttt{Mathlib.RingTheory.Noetherian} for the
    Noetherian-induction termination argument.
  \item \texttt{AlgebraicGeometry.IsProjective} (or
    \texttt{ProjectiveMorphism}) and the relative-\(\PP^n\) construction, to
    state and unpack the projective-morphism hypothesis on
    \(\pi : \PP^n_S \to S\).
  \item Higher-direct-image base-change lemmas (\texttt{Mathlib.AlgebraicGeometry
    .Cohomology}, or a Stacks-tag-02KH-shaped lemma; cf.\ the existing
    \cref{chap:Picard_RelativeSpec} use of the same tag) for the \(E_i = \pi_*\F(N+i)\)
    construction in \cref{lem:noetherian_induction_strata}.
  \item Castelnuovo--Mumford regularity / Serre vanishing for the
    \texttt{N\_1}-existence step.  Mathlib does not yet have these in the
    relative form needed; a future sub-lemma may stub them or rely on
    \texttt{Module.IsTorsion} arguments on the graded module
    \(\bigoplus_{m \ge N} \pi_* \F(m)\).
\end{itemize}

\section{Out of scope}
\label{sec:flatstrat_out_of_scope}

This chapter packages \emph{only} the flattening stratification of a coherent
sheaf on a projective morphism, plus the relative-curve specialisation. The
following downstream / parallel constructions live in separate sibling
chapters and are explicitly out of scope here:
\begin{itemize}
  \item \textbf{Castelnuovo--Mumford regularity} (Nitsure~\S2) --- input to the
    Quot construction, used implicitly inside the proof of
    \cref{lem:noetherian_induction_strata} as the assertion that the
    Hilbert function stabilises beyond a uniform \(N\).  Its own blueprint home
    is the planned chapter \texttt{Picard\_CastelnuovoMumford.tex} (not yet
    scaffolded).
  \item \textbf{Quot scheme construction} (Nitsure~\S5, A.2.b) --- consumer of
    the present chapter; its home is the planned \texttt{Picard\_QuotScheme.tex}.
  \item \textbf{Relative Picard functor} (A.1.c) ---
    \texttt{Picard\_RelPicFunctor.tex} (planned).
  \item \textbf{Picard scheme representability for a smooth proper curve}
    (A.2.c) --- the assembly theorem fed by both the Quot construction (which
    consumes this chapter) and the relative Picard functor; home is the
    planned \texttt{Picard\_FGAPicRepresentability.tex}.
\end{itemize}
